\section{Experimental Evaluation}
\label{sec:eval}
The evaluation asks whether registered-harness tasks and direct operation calls
can execute bounded recipe actions across two simulator scenes with different
signal and protocol mappings while retaining route-appropriate records. R20 is
the primary result: six sequential tasks through one registered harness, three
per scene on one simulator instance per scene, with task-to-approval-to-receipt-
to-sample links. R21 is a deterministic fixture replay of the registered-
harness path; R22 separately exercises the direct operation endpoint without a
task or harness. All runs use repository simulators. Supplementary File~S1
contains their separate records, scripts, and trace checks.

\subsection{Evidence Units and Experimental Procedure}
\label{sec:testbed}
\label{sec:procedure}
R20 uses the registered harness adapter named \texttt{omp} in cast-film and
annealing scenes
(\texttt{omp} is the project adapter key, not a network protocol),
with three sequential task executions per scene on one simulator instance per
scene. Each execution creates a new logical line, product, recipe, active run,
channel, agent, and task on that shared scene simulator. Line start reapplies the recipe baseline
(120~rpm or 730~$^{\circ}$C), so the control read returns to baseline while
the simulated process state carries forward. These are sequential software
runs, not independent process replications. R21 and R22 use fresh isolated
application/simulator environments and are reported separately from R20.

Table~\ref{tab:routes} summarizes the evaluated paths and their inference
limits; Table~\ref{tab:q8} gives the six R20 process snapshots.

\begin{table*}[!t]
\centering
\caption{Evaluation paths, recorded outcomes, and inference limits. The route
counts are kept separate because the runs use different entry paths and
execution modes.}
\label{tab:routes}
\footnotesize
\setlength{\tabcolsep}{3pt}
\begin{tabularx}{\textwidth}{@{}p{1.25cm}p{3.4cm}X p{4.0cm}@{}}
\toprule
Path & Execution & Recorded outcome & Inference limit \\
\midrule
R20 & Six sequential tasks through one registered harness; three per scene. & 6/6 task--approval--receipt--sample links; six verified receipts; all six single probes were outside target. & One simulator per scene; one probe per task; no process-attainment claim. \\
R21 & One deterministic fixture task per scene on the registered-harness route. & 2/2 fixture traces; 31 runner assertions passed. & No external model decision or independent process replication. \\
R22 & One direct operation call per scene, without a task, harness, or language-model decision. & 2/2 calls approved; each returned a 1/1 verified receipt and matching readback, then one active-run-tagged snapshot; both were outside target; 19 checks passed. & No task-to-sample link; one snapshot per scene; no alternate-harness portability claim. \\
Context tests & Twelve stale-context mutations on three recipe routes; production handler and test-only ablation. & Production handler refused 12/12 before writes; one queued post-recheck change was also refused. Without tuple equality, six product/run mutations triggered 12 mock writes and six execution records. & Temporary/mock state; component ablation, not a platform baseline or physical safety test. \\
\bottomrule
\end{tabularx}
\end{table*}

\subsection{Registered-Harness Process Path (R20)}
\label{sec:harness-e2e}
R20 ran the local v0.8.0 working-tree application with the registered
\texttt{omp} adapter and GLM-5.3-Flash. Supplementary File~S1 lists the source
reconstruction requirements; credentials are omitted and external model
decisions are not deterministic. Cast-film used OPC UA control and MQTT
thickness sensing; annealing used Modbus TCP for control and sensing. In each
scene the model selected one value from a prevalidated candidate set for a
bounded \texttt{recipe\_trial}, and a synthetic administrator approved it.
Each route record contains one 1/1 readback-verified receipt. After verifying
the receipt, the runner waited 20~s, probed and ingested one sample, and queried
it by active-run ID through the isolated SQLite fallback. All six tasks
completed and passed 76 runner assertions. The trace audit found 6/6 complete
task--approval--receipt--sample links. The trial-operation run ID differs from
the active production-run ID used for the sample; both link through the task
and approval. Task state was checked after sampling, so completion-to-sample
latency is unavailable. Table~\ref{tab:q8} reports all six values and
identifiers are in the S1 trace matrix. All samples remained outside target;
one probe per task does not measure process response or normal online DAQ.

\begin{table}[!t]
\centering
\caption{Bounded recipe-trial settings and post-readback simulator probes. All
six samples remained outside their target bands.}
\label{tab:q8}
\footnotesize
\setlength{\tabcolsep}{2.5pt}
\begin{tabularx}{\columnwidth}{@{}lXcX@{}}
\toprule
Case & Pre-task value / target band & Control: start$\to$trial & Sample \\
\midrule
Film-1 & $55.82/48$--$52~\mu$m & $120\to125$~rpm & $45.69~\mu$m \\
Film-2 & $45.69/48$--$52~\mu$m & $120\to130$~rpm & $46.69~\mu$m \\
Film-3 & $46.69/48$--$52~\mu$m & $120\to125$~rpm & $45.69~\mu$m \\
Anneal-1 & $675.40/700$--$760~^{\circ}$C & $730\to745$~$^{\circ}$C & $679.60~^{\circ}$C \\
Anneal-2 & $678.30/700$--$760~^{\circ}$C & $730\to745$~$^{\circ}$C & $677.70~^{\circ}$C \\
Anneal-3 & $676.90/700$--$760~^{\circ}$C & $730\to745$~$^{\circ}$C & $678.10~^{\circ}$C \\
\bottomrule
\end{tabularx}
\end{table}

\subsection{Deterministic Fixture Replay (R21)}
R21 used a deterministic transport fixture for the registered \texttt{omp}
adapter, without an LLM decision. One task per scene exercised approval,
verified receipt/readback, and a run-tagged sample; all 31 runner checks,
including cleanup, passed. This is a replay of the same harness route, not an
independent agent or process replication, and is reported separately from R20.

\subsection{Direct Operation Endpoint (R22)}
R22 used a fresh isolated environment. A worker identity issued one direct
\texttt{recipe\_trial} call per scene through
\texttt{/api/workshop/agent-tools/invoke}, without a task, harness transport,
or external model. Synthetic approval produced a 1/1 verified receipt and
matching readback for each call. After 20~s, one sample per scene was queried
by active-run ID: 46.37~$\mu$m versus a 48--52~$\mu$m band and 675.20~$^{\circ}$C
versus a 700--760~$^{\circ}$C band. All 19 checks passed. This tests a direct
entry route and two mappings, not process response or cross-harness
portability; R22 is separate from R20 and R21.

\subsection{Software Contract Checks}
The local unit suite passed 566/566 tests across seven suites. In the 20-case
approval suite, twelve injected line, product, recipe, or run changes across
three recipe routes were refused before mock writes or new execution records.
One queued post-recheck context change was also refused before dispatch. In a
test-only ablation that removed tuple equality but retained authority and
limit checks, six product/run changes caused twelve mock writes and six
execution records; independent checks still blocked line/recipe changes.
This component ablation shows the tuple check's incremental software-state
coverage, not a platform baseline or safety guarantee. Tests use mock or
temporary state and do not establish atomicity during device I/O or test
end-to-end DAQ tag failures.
